\section{From order flow to a Gaussian price: the volume bridge}
\label{sec:bridge}

\subsection{The price as accumulated order flow}

Over $[t,t+\Delta]$ let the market execute $M$ trades, trade $k$ carrying signed
volume $\epsilon_k v_k$ with $\epsilon_k\in\{-1,+1\}$ the sign (buyer- or
seller-initiated) and $v_k>0$ the size. Write $V=\sum_k v_k$ for total volume
and $Q=\sum_k \epsilon_k v_k$ for net signed order flow.

\begin{assumption}[Linear impact]\label{as:impact}
The log-price responds linearly to net signed order flow,
$\Delta X_{t,\Delta} = X_{t+\Delta}-X_t = \lambda\,Q$, with $\lambda>0$ constant
over the interval.
\end{assumption}

This is Kyle's \citeyearpar{kyle1985} equilibrium: the market maker, unable to
distinguish informed from uninformed flow, sets a linear pricing rule and
$\lambda$ is the inverse depth. It is the definition of ``price as a function of
traded volume'' that the rest of the section needs, and it is the weakest link
in the chain; Section~\ref{sec:bridge-caveats} says what happens when it fails.

\begin{assumption}[Exchangeable signs and a volume-proportional variance]
\label{as:signs}
Conditional on $V$ and on the information available at $t$, the market's belief
about $Q$ is constrained only by its first two moments, with
$\mathbb{E}[Q\mid V] = m$ and $\mathrm{Var}(Q\mid V) = \varsigma^2 V$.
\end{assumption}

The variance form is what one gets when trade signs are exchangeable with
bounded correlation and sizes have a stable mean: for independent signs,
$\mathrm{Var}(Q)=\sum_k v_k^2 \approx \bar v V$. The economic content is that
\emph{volume, not time, is the natural measure of how much can happen}.

\subsection{Maximum entropy in volume time}

\begin{proposition}[Maximum entropy under first and second moments]
\label{prop:maxent}
Among absolutely continuous distributions on $\mathbb{R}^N$ with mean $\mu$ and
covariance $\Sigma\succ0$, the Gaussian $\mathcal N(\mu,\Sigma)$ uniquely
maximises differential entropy, and
$h^{\max}(\Sigma)=\tfrac12\log\left((2\pi e)^N\det\Sigma\right)$.
\end{proposition}

\begin{proof}
Let $q=\mathcal N(\mu,\Sigma)$ and let $p$ have the same first two moments and
finite differential entropy. Since $\log q$ is an exact quadratic, its
expectation depends only on those moments, so
$\mathbb{E}_p[\log q]=\mathbb{E}_q[\log q]=-h(q)$. Non-negativity of relative
entropy gives
$0\le \KL{p}{q} = -h(p)+h(q)$, hence $h(p)\le h(q)$, with equality iff $p=q$
almost everywhere. The argument requires no differentiability of $p$ and no
prior existence of a maximiser.
\end{proof}

The variational route -- a Lagrangian in the normalisation, mean and covariance
constraints, whose first-order condition makes $\log p$ quadratic -- explains
\emph{why} the maximiser is Gaussian, but the divergence argument is what
establishes optimality and uniqueness.

Applying Proposition~\ref{prop:maxent} to Assumption~\ref{as:signs} in one
dimension and mapping through Assumption~\ref{as:impact}:

\begin{equation}
\label{eq:mdh}
\Delta X_{t,\Delta}\ \big|\ V \ \sim\ \mathcal{N}\!\left(\lambda m,\ \lambda^2\varsigma^2 V\right).
\end{equation}

Equation~\eqref{eq:mdh} is the bridge. Three things about it deserve emphasis.

First, \emph{the conditioning variable is volume, not time}. Maximum entropy
does not deliver a variance; it delivers the least-committal distribution given
one. Assumption~\ref{as:signs} supplies the variance from an observable, and
that observable is traded volume. Equation~\eqref{eq:mdh} is the
mixture-of-distributions hypothesis of \citet{clark1973} and
\citet{tauchen1983}, obtained here as a maximum-entropy statement rather than
posited; \citet{ane2000} show empirically that returns measured in trade time
are close to Gaussian even though calendar-time returns are not, which is the
prediction \eqref{eq:mdh} makes.

Second, \emph{Gaussianity is conditional}. Unconditionally,
$\Delta X$ is a scale mixture of normals over the random volume $V$, hence
leptokurtic even when \eqref{eq:mdh} holds exactly. This is the
\citet{mandelbrot1967} subordination mechanism, and it means fat tails are not
by themselves evidence against maximum entropy.

Third, the multivariate extension is immediate: with $N$ assets,
$\Delta \mathbf{X}\mid V \sim \mathcal N(\boldsymbol\mu V, \Sigma V)$, and
Proposition~\ref{prop:maxent} applies as stated.

\subsection{From Gaussian increments to Brownian motion and to GBM}

Let $\theta_t=\int_0^t v_s\,ds$ be the cumulative volume clock. Equation
\eqref{eq:mdh} gives the one-increment distribution; a process needs more.

\begin{assumption}[Temporal consistency]\label{as:time}
In the volume clock, increments over disjoint intervals are independent, sample
paths are continuous, and the covariance of an increment is proportional to
elapsed volume.
\end{assumption}

Under Assumption~\ref{as:time} with constant coefficients there is a matrix $B$
with $BB^{\!\top}=\Sigma$ such that $X_\theta = X_0 + \mu\theta + BW_\theta$ in
volume time, so in calendar time
\begin{equation}
\label{eq:subordination}
X_t = X_0 + \mu\,\theta_t + B\,W_{\theta_t},
\end{equation}
a Brownian motion \emph{subordinated} to the volume clock. If trading intensity
is constant, $\theta_t=\bar v\,t$ and \eqref{eq:subordination} is ordinary
Brownian motion in calendar time; exponentiating and applying Itô,
\begin{equation}
\frac{dS_{i,t}}{S_{i,t}} = \left(a_{i,t}+\tfrac12\sigma_{i,t}^2\right)dt + \sigma_{i,t}\,dW_{i,t},
\end{equation}
so the drift of the price is not the drift of the log-price -- the same
second-order correction that will reappear in Section~\ref{sec:channels} when we
differentiate $\log\det R_t$.

It is worth being precise about what Assumption~\ref{as:time} does. By Monroe's
\citeyearpar{monroe1978} theorem, \emph{every} continuous local martingale is a
time-changed Brownian motion, so subordination alone has no empirical content.
The content is in the clock: \eqref{eq:mdh} says the right clock is volume, and
Assumption~\ref{as:time} says the market has no memory in that clock. Maximum
entropy is what selects Brownian motion from the class Monroe's theorem allows.

\subsection{The contrapositive}
\label{sec:contrapositive}

Proposition~\ref{prop:maxent} runs both ways. Write $q_t$ for the Gaussian with
the conditional mean and covariance of $p_t$. Then

\begin{equation}
\label{eq:deficit}
\DD_t \;\equiv\; h(q_t) - h(p_t) \;=\; \KL{p_t}{q_t} \;\ge\; 0,
\end{equation}
with equality if and only if $p_t$ is Gaussian. So: \emph{prices are non-Gaussian
exactly to the extent that their conditional distribution is not
entropy-maximal}, and $\DD_t$ measures the extent. In the language of
Section~\ref{sec:bridge}, a positive $\DD_t$ means constraints beyond the first
two moments are binding -- that somebody knows something the second moment does
not encode.

Two sources produce a positive deficit, and separating them matters. A random
volume clock produces a scale mixture, which is symmetric and leptokurtic.
Binding directional information produces asymmetry. Section~\ref{sec:channels}
splits $\DD$ along exactly this line.

\subsection{Where the bridge is weak}
\label{sec:bridge-caveats}

Assumption~\ref{as:impact} is linear; the empirical impact of a metaorder is
closer to square-root in size, which makes the mapping from flow to price
concave and the induced return distribution thinner-tailed than the flow.
Assumption~\ref{as:signs} treats signs as exchangeable; order flow is in fact
strongly autocorrelated, which inflates $\mathrm{Var}(Q)$ above $\varsigma^2 V$
and makes the volume clock an imperfect variance proxy. Neither failure changes
the \emph{form} of the argument -- maximum entropy given a variance budget is
still Gaussian -- but both mean the volume clock is an approximation to the
information clock rather than the thing itself. We use volume-time reasoning to
motivate the framework and calendar-time estimation throughout the empirical
work, so no result below depends on \eqref{eq:mdh} holding exactly.

\section{Three channels of entropy}
\label{sec:channels}

\subsection{Scale and dependence}

Write $\Sigma_t = D_tR_tD_t$ with $D_t=\mathrm{diag}(\sigma_{1,t},\dots,\sigma_{N,t})$.
Then $\det\Sigma_t = \left(\prod_i\sigma_{i,t}^2\right)\det R_t$ and
\begin{equation}
\label{eq:ceiling}
h(q_t) \;=\; \underbrace{\tfrac{N}{2}\log(2\pi e)}_{\text{constant}}
\;+\; \underbrace{\sum_i \log\sigma_{i,t}}_{H^{vol}_t}
\;+\; \underbrace{\tfrac12\log\det R_t}_{H^{dep}_t}.
\end{equation}
By Hadamard's inequality $H^{dep}_t\le 0$, with equality iff $R_t=I$:
\emph{correlation always destroys entropy}, and $-H^{dep}_t$ is precisely the
Gaussian multi-information (total correlation) of the return vector. A useful
reading of the same quantity is spectral: with $\lambda_k$ the eigenvalues of
$R_t$ and $\tilde\lambda_k=\lambda_k/N$, the participation number
$N^{\text{eff}}_t=\exp\!\left(-\sum_k\tilde\lambda_k\log\tilde\lambda_k\right)$
is the effective number of independent risk modes; it equals $N$ for $R=I$ and
falls toward $1$ as one collective mode absorbs the variance.

Equation~\eqref{eq:ceiling} is the ceiling, not the entropy. It is what the
covariance matrix alone implies, and Proposition~\ref{prop:maxent} says the
truth sits below it.

\subsection{Shape: the entropy deficit and what it is made of}

Subtracting \eqref{eq:deficit} from \eqref{eq:ceiling} gives the paper's
decomposition, three channels and a constant:
\begin{equation}
\label{eq:three}
h(p_t) \;=\; \tfrac{N}{2}\log(2\pi e) \;+\; \underbrace{H^{vol}_t}_{\text{scale}} \;+\; \underbrace{H^{dep}_t}_{\text{dependence}} \;-\; \underbrace{\DD_t}_{\text{shape}}.
\end{equation}

The shape channel is skewness and fat tails, and the Edgeworth expansion says so
explicitly: for a standardised variable with skewness $S$ and excess kurtosis
$K$ near Gaussianity,
\begin{equation}
\label{eq:edgeworth}
J \;\approx\; \frac{S^2}{12} + \frac{K^2}{48}.
\end{equation}
Equation~\eqref{eq:edgeworth} is the answer to the question ``where would an
entropy jump be visible?'': in the third and fourth moments, additively, with an
asymmetry channel and a tail channel. It is also the source of the sign problem
that Section~\ref{sec:dynamics} has to solve, because $S$ enters squared --
\emph{negentropy cannot tell a crash from a melt-up}. As an estimator
\eqref{eq:edgeworth} is unusable on daily returns; Section~\ref{sec:program}
documents how badly and what we use instead.

\subsection{The structural index}

The scale channel is a nuisance for our purposes: it dominates the level of
\eqref{eq:three} numerically, it moves with volatility, and a change in units
changes it. The other two channels are scale-free. Combining them, and using the
exact chain rule
$\KL{p}{\prod_i q_i} = \sum_i \KL{p_i}{q_i} + I(x)$ with $I$ the
multi-information:

\begin{definition}[Structural information index]
\label{def:J}
\begin{equation}
\JJ_t \;\equiv\; \DD_t \;-\; H^{dep}_t \;=\; \KL{p_t}{\textstyle\prod_i q_{i,t}} \;\ge\; 0 ,
\end{equation}
the divergence of the conditional joint return distribution from a product of
Gaussian marginals with the same scales.
\end{definition}

$\JJ$ is zero if and only if returns are independent and Gaussian; it rises when
assets couple, when marginals fatten, or when they skew. It is measured in nats,
it is invariant to the level of volatility, and it is additive across its two
channels, so one can always ask how much of a move came from dependence and how
much from shape. It is the state variable of Section~\ref{sec:dynamics}.

